\documentclass[12pt]{article}
% Préambule commun à tous les documents extraits de « La Revue CAPES »
\usepackage[T1]{fontenc}
\usepackage[utf8]{inputenc}
\usepackage{lmodern}
\usepackage{amsmath,amssymb,amsthm,amsfonts,mathrsfs}
\usepackage{setspace}
\usepackage{listings}
\usepackage{pgfplots}
\pgfplotsset{compat=1.18}
\usepackage{longtable}
\usepackage{adjustbox}
\usepackage{lettrine}
\usepackage{tkz-tab}
\usepackage{changepage}
\usepackage{pdflscape}
\usepackage{graphicx}
\usepackage{tikz}
\usepackage{xcolor}
\usepackage[a4paper,top=2cm,bottom=2.5cm,left=2.5cm,right=2cm]{geometry}
\usepackage{fancyhdr}
\usepackage{draftwatermark}
\SetWatermarkText{La RevueCapes}
\SetWatermarkLightness{0.9}
\SetWatermarkScale{2}
\usepackage{hyperref}
\hypersetup{colorlinks=true,linkcolor=blue!50!black,urlcolor=blue!50!black}

\definecolor{revue}{RGB}{20,60,120}

\pagestyle{fancy}
\fancyhf{}
\renewcommand{\headrulewidth}{0.4pt}
\setlength{\headheight}{15pt}

% \entete{Matière}{Titre}{Type de document}
\newcommand{\entete}[3]{%
  \fancyhead[L]{\small\textsc{La Revue CAPES} -- #1}%
  \fancyhead[R]{\small #2}%
  \fancyfoot[C]{\thepage}%
  \thispagestyle{plain}%
  \begin{center}
    {\color{revue}\rule{\textwidth}{1.2pt}}\\[4pt]
    {\small\textsc{Concours directs CAP-CEG / CAPES -- Burkina Faso}}\\[6pt]
    {\LARGE\bfseries\color{revue} #2}\\[6pt]
    {\large\itshape #1 -- #3}\\[2pt]
    {\color{revue}\rule{\textwidth}{1.2pt}}
  \end{center}
  \vspace{0.5cm}%
}

% Encadré pour les remarques ajoutées dans les corrigés
\newcommand{\remarque}[1]{\par\medskip\noindent\fcolorbox{revue}{revue!6}{\parbox{\dimexpr\linewidth-2\fboxsep-2\fboxrule}{\textbf{\color{revue}Remarque.} #1}}\par\medskip}


\begin{document}
\entete{Mathématiques}{Sujet type n°8}{Énoncé}
\begin{spacing}{1.5}

\medskip\textbf{Exercice 1}

Soient $e_1, e_2, e_3$ trois vecteurs de $\mathbb{R}^3$. On suppose que $B = (e_1,e_2,e_3)$ est une base de $\mathbb{R}^3$.
 
 On note $T$ l'application linéaire définie par $T(e_1) = T(e_3) = e_3$ et $T(e_2) = -e_1+e_2+e_3$.

1. (a) Décrire le noyau de $T$.
 
 (b) Donner la matrice de $T$ dans la base $B$. On la note $A$.
 
 2. On pose $f_1 = e_1-e_3, f_2 = e_1-e_2, f_3 = -e_1 +e_2 +e_3$.
 
 (a) Calculer $e_1, e_2, e_3$ en fonction de $f_1, f_2, f_3$. Les vecteurs $f_1, f_2, f_3$ forment-ils
 une base de $\mathbb{R}^3$ ?
 
 (b) Écrire la matrice de $T$ dans la base $(f_1,f_2,f_3)$. On la note $B$.
 
3. On note $P= \begin{pmatrix}
1 & 1 & -1 \\ 
0 & -1 & 1 \\ 
-1 & 0 & 1
\end{pmatrix}$

Vérifier que $P$ est inversible et calculer $P^{-1}$. Quelle relation relie $A,B,P$ et $P^{-1}$?

\medskip\textbf{Exercice 2}

Le plan $(P)$ est rapporté à un repère orthonormé direct $(\overrightarrow{0} ; \overrightarrow{u}; \overrightarrow{v})$. 
on désigne par $T$ l'application de $(P)$ dans $(P)$ qui à tout point $M$ d'affixe $z $associe le point $M'$ d'affixe $z' = (1 + i)z-i$. 

1. Montrer que$ T$ est la composée d'une rotation et d'une homothétie de rapport positif dont on donnera les éléments caractéristiques. On note $\Omega$ le point invariant par $T$. 

Donner une mesure de l'angle orienté $(\overrightarrow{\Omega M},\overrightarrow{MM'}) \text{ en supposant que }  M\neq \Omega$.


2.a) Construire $M'$ image de $M$ par $T$ où $M$ est un point donné distinct de $\Omega$. 

b) Déterminer l'image $(D')$ par $T$ de la droite $(D)$ d'équation $y = x$. Construire $(D')$. 

3.a) Montrer qu'il existe un point $B$ de $(P)$ distinct de $\Omega$ et un seul 
tel que les affixes $z_0$ de $B$ et $z'_0$ de $B'$ (où $B'$ est l'image de $B $par $T$) soient liés par la relation $z_0z'_0 = 1$. Placer les points $B$ et $B'$ dans le repère $(\overrightarrow{0} ; \overrightarrow{u}; \overrightarrow{v})$. 

b) Soit $\Omega'$ le symétrique de $\Omega$, par rapport à O. Montrer que les points $\Omega',\Omega,B et B' $sont cocycliques. 
Déterminer l'affixe du centre G du cercle $(\Gamma)$ passant par $\Omega', \Omega, B et B'$. 

\medskip\textbf{Exercice 3 }

1. Calculer

$I=\int\int_D\frac{dxdy}{(1+x^2+y^2)^2}$ ou $D=\lbrace {(x,y) \in \mathbb{R}^2,\vert x\vert \leq x^2+y^2\leq 1\rbrace}$.

2. Calculer l'intégrale double

$J=\int\int_D(x+y)^2dxdy$ ou $D=\lbrace (x,y)\in \mathbb{R}^2, x^2+y^2-x \leq 0;x^2+y^2-y \geq 0$ et $y \geq 0 \rbrace$.

3. Soit $D=\lbrace (x,y,z)\in \mathbb{R}^3;x\geq 0,y\geq 0,z\geq 0,x+y+z\leq 1 \rbrace$.

Utilisez le changement de variable $u=x+y+z; v=\frac{z}{y+z}; w=\frac{y+z}{x+y+z}$ pour calculer

$S=\int\int_D\frac{z^3}{(x+y+z)(y+z)}dxdydz$.

4. Calculer l'aire du domaine elliptique :

$E = \lbrace (x,y)\in \mathbb{R}^2;\frac{x^2}{a^2}+\frac{y^2}{b^2}\leq 1\rbrace$.


\medskip\textbf{Exercice 4 }

On pose $f(n)=\int_{0}^{1}x^ne^{-x}dx \hspace{0.05cm} \forall n\in \mathbb{N}$.

1. Montrer que la suite $(f(n))$ est positive et décroissante. Au moyen d'une intégration par parties
 donner une relation de récurrence entre $f(n)$ et $f(n-1)$.
 
 2. Montrer par récurrence que pour tout $n$,  $f(n)= \frac{n!}{e}(e-\sum_{k=0}^{n}\frac{1}{k!})$
 
 3. Montrer que l'on a $\frac{1}{e(n+1)}\leq f(n)\leq \frac{1}{n+1}$.
 
 4. En déduire la nature des série $\sum_{n=1}^{+\infty}f(n), \sum_{n=1}^{+\infty}\frac{f(n)}{n}$ et $\sum_{n=1}^{+\infty}(-1)^nf(n)$.
 
 5. Déterminer le rayon de convergence de la série entière $\sum_{n=1}^{+\infty}f(n)x^n$.
 
 
 \quad
 \quad
 

\end{spacing}
\end{document}
